\documentclass[../../hsm_tok_q.tex]{subfiles}

\begin{document}
    \setcounter{figure}{0}
    \textsf{図\ref{fig/hsm_tok_2011_1st_03_01_q}}で，
    点Oは原点，点Aの座標は$(0,\;-4)$であり，
    直線$\ell$は一次関数$y=\myfraca{1}{2}x+3$のグラフを表している．
    
    直線$\ell$と$x$軸との交点をBとする．
    
    直線$\ell$上にある点をPとし，2点A，Pを通る直線を$m$とする．
    
    次の各問に答えよ．
    %
    \vspace{.5\baselineskip}
    {\par \centering
        \setlength{\abovecaptionskip}{5pt}
        \includegraphics[width=.8\linewidth]%
            {\subfix{fig_hsm_tok_2011_1st_03/fig_hsm_tok_2011_1st_03_01_q.pdf}}
            \captionof{figure}{} \label{fig/hsm_tok_2011_1st_03_01_q}
    \par}
    %
    \begin{enumerate}
        \item 点Pの$y$座標が$-1$のとき，
        点Pの$x$座標を，
        次の\textsf{ア}～\textsf{エ}のうちから選び，記号で答えよ．
        
        \vspace{.15\baselineskip}
        \begin{tabularx}{\dimexpr\linewidth-3.5mm}{@{}LLLL@{}}
            {\textsf ア}\quad$-8$ & {\textsf イ}\quad$-\myfraca{9}{2}$&
            {\textsf ウ}\quad$-2$ & {\textsf エ}\quad$\myfraca{5}{2}$
        \end{tabularx}

        \item 点Pが点Bに一致するとき，直線$m$の式を，
        次の\textsf{ア}～\textsf{エ}のうちから選び，記号で答えよ．
        
        \vspace{.15\baselineskip}
        \begin{tabularx}{\dimexpr\linewidth-3.5mm}{@{}LL@{}}
            {\textsf ア}\quad$y=-\myfraca{3}{2}x-4$ & {\textsf イ}\quad$y=-\myfraca{3}{2}x-6$ \\[8pt]
            {\textsf ウ}\quad$y=-\myfraca{2}{3}x-4$ & {\textsf エ}\quad$y=-\myfraca{2}{3}x-6$
        \end{tabularx}

        \newpage

        \item \textsf{図\ref{fig/hsm_tok_2011_1st_03_02_q}}は，
            \textsf{図\ref{fig/hsm_tok_2011_1st_03_01_q}}において，
            点Pの$x$座標が正の数のとき，
            $x$軸上にあり$x$座標が点Pの$x$座標と等しい点をQとし，
            点Aと点B，点Aと点Q，点Pと点Qをそれぞれ結んだ場合を表している．
            
            $\mathrm{△APB}$の面積が$\mathrm{△AQP}$の面積の2倍になるとき，点Pの$x$座標を求めよ．
            %
            \vspace{.5\baselineskip}
            {\par \centering
                \setlength{\abovecaptionskip}{5pt}
                \includegraphics[width=.8\linewidth]%
                    {\subfix{fig_hsm_tok_2011_1st_03/fig_hsm_tok_2011_1st_03_02_q.pdf}}
                    \captionof{figure}{} \label{fig/hsm_tok_2011_1st_03_02_q}
            \par}
        \end{enumerate}

\end{document}

